\documentclass[11pt, a4paper]{article}

% --- Packages ---
\usepackage[utf8]{inputenc}
\usepackage[T1]{fontenc}
\usepackage{lmodern}
\usepackage[french]{babel}
\usepackage[left=1.5cm, right=1.5cm, top=1.5cm, bottom=1.5cm]{geometry}
\usepackage{hyperref}     % Pour les liens cliquables
\usepackage{titlesec}     % Pour personnaliser les titres de section
\usepackage{enumitem}     % Pour personnaliser les listes
\usepackage{parskip}      % Pour espacer les paragraphes
\usepackage{xcolor}
\usepackage{fontawesome}


% --- Configuration des couleurs et liens ---
\definecolor{linkcolor}{rgb}{0.0, 0.0, 0.5} % Bleu sombre et pro pour les liens
\hypersetup{
    colorlinks=false,
    hidelinks
}


% --- Configuration des titres de sections ---
\titleformat{\section}{\large\scshape\raggedright}{}{0em}{}[\titlerule]
\titlespacing{\section}{0pt}{12pt}{6pt}



\newcommand{\cventry}[4]{
    \vspace{0.1cm}
    \noindent \textbf{#1} \hfill #2 \\
    \noindent \textit{#3} \hfill \textit{#4} \par
    \vspace{0.05cm}
}


\setlist[itemize]{label=\textbullet}

\begin{document}
\pagestyle{empty} % Pas de numéro de page

% ================= EN-TÊTE =================


\begin{center}
    {\Huge \textbf{Marcel ASSIE}} \\[0.2cm]
    {\Large \textsc{Ingénieur en Géomatique}} \\[0.3cm]
    \footnotesize
    \faMobile~+33 7 68 06 25 78
    ~$|$~
    \faEnvelope~\href{mailto:marcelassie2k@gmail.com}{marcelassie2k@gmail.com}
    ~$|$~
    \faLinkedinSquare~\href{https://www.linkedin.com/in/marcel-assie/}{LinkedIn}
    ~$|$~
    \faGithub~\href{https://github.com/MarcelAssie/}{GitHub}
    ~$|$~
    \faGlobe~\href{https://marcel-assie-portefolio.dev}{Portfolio}
    ~$|$~
    \faMapMarker~Champs-sur-Marne (77)
\end{center}

\vspace{0.3cm}
\noindent Ingénieur géomaticien spécialisé en Geo Data Science, passionné par l'écosystème Geo-Data et les technologies géospatiales. À la recherche d'une opportunité en \textbf{CDI} pour mettre à profit mes compétences en \textbf{développement géospatial}, \textbf{automatisation de pipelines de données} et \textbf{analyse spatiale}, au service de projets à fort impact permettant de transformer des données territoriales complexes en véritables outils décisionnels.

% ================= FORMATIONS =================
\section*{Formations}

\cventry{École Nationale des Sciences Géographiques (Geodata Paris)}{Champs-sur-Marne, France}{Ingénieur en Géomatique \textbullet\ Spécialité Geo Data Science}{Sept. 2023 -- Sept. 2026}
\cventry{Institut National Polytechnique Félix Houphouët-Boigny}{Yamoussoukro, Côte d'Ivoire}{Licence de Géomètre}{Oct. 2019 -- Juil. 2022}
\cventry{Lycée Scientifique de Yamoussoukro}{Yamoussoukro, Côte d'Ivoire}{Baccalauréat, Série Scientifique}{Sept. 2016 -- Juil. 2019}

% ================= EXPÉRIENCES PROFESSIONNELLES =================
\section*{Expériences Professionnelles}

\cventry{ISOGEO}{Paris, France}{Ingénieur Développeur SIG -- Alternance}{Sept. 2025 -- Sept. 2026}
\begin{itemize}[leftmargin=0.5cm, nosep]
    \item Refonte technologique du « Scan Isogeo » en remplaçant l'ETL FME par de nouvelles chaînes de traitements automatisées.
    \item Manipulation et traitement d'une grande variété de formats de données géographiques (vecteur, raster, bases de données) avec GDAL et SQL.
    \item Exploration et implémentation de modèles d'IA pour extraire automatiquement et enrichir les métadonnées sémantiques.
    \item Rédaction de spécifications fonctionnelles et techniques, et suivi du projet au sein d'une équipe fonctionnant en méthode Agile.
\end{itemize}

\cventry{ORYJIN}{Lille, France}{Ingénieur Geo Data Scientist -- Stage}{Avr. -- Août 2025}
\begin{itemize}[leftmargin=0.5cm, nosep]
    \item Recherche, audit qualité et intégration de diverses sources de données géospatiales et socio-démographiques (INSEE, IGN) dans le référentiel mondial open source H3.
    \item Développement de pipelines ETL et mise au point d'algorithmes statistiques pour calculer des agrégats géographiques (Python, Snowflake, SQL).
    \item Extraction d'informations spatiales (zones d'habitation, infrastructures) à partir d'images satellites à l'aide de modèles d'IA de Computer Vision.
    \item Automatisation de la datavisualisation via l'IA générative : création de modèles descriptifs et de flux de prompts pour générer des images réalistes décrivant des territoires.
\end{itemize}

\cventry{IGN}{Forcalquier, France}{Ingénieur Géomaticien -- Stage}{Mai -- Juil. 2024}
\begin{itemize}[leftmargin=0.5cm, nosep]
    \item Mini-projets appliqués en topométrie, photogrammétrie, métrologie (auscultation de ponts), télédétection et géodésie, avec formalisation méthodologique et analyse des incertitudes.
    \item Conception d’un workflow automatisé de cartographie des restrictions de circulation (Alpes-de-Haute-Provence) :
    \begin{itemize}[leftmargin=0.5cm, nosep]
        \item Structuration et nettoyage de jeux de données multisources.
        \item Modélisation des traitements géospatiaux (QGIS Model Builder).
        \item Développement d’algorithmes de contrôle qualité et de validation des données.
    \end{itemize}
\end{itemize}

% ================= COMPÉTENCES =================
\section*{Compétences}
\noindent \textbf{Langages :} Python, SQL, JavaScript, HTML/CSS, R, Cypher \\
\textbf{Technologies :} Docker, GitHub, PostgreSQL/PostGIS, Django, Flask, Vue.js, React Native, Neo4j \\
\textbf{Bibliothèques :} GDAL, Proj, GeoPandas, ArcPy, PyQGIS, Shapely, Leaflet.js \\
\textbf{Outils SIG :} QGIS, ArcGIS Pro, FME, GeoServer, AutoCAD \\
\textbf{Aptitudes professionnelles :} Forte capacité d’adaptation, Rigueur, Esprit critique et d'analyse, Prise d’initiative

% ================= PROJETS =================
\section*{Projets}

\begin{itemize}[leftmargin=0.5cm, nosep]
    \item \textbf{Extraction de données contextuelles (ODD)} :
    \begin{itemize}[leftmargin=0.5cm, nosep]
        \item Développement d'une méthode avancée d'extraction basée sur les Grands Modèles de Langage (LLM) et structuration sous forme de graphes de connaissances pour analyser les relations entre les ODD (PyTorch, Transformers, NLP, LLM, Neo4j).
    \end{itemize}
    \vspace{0.2cm}
    \item \textbf{Application mobile de navigation (Gendarmerie nationale)} :
    \begin{itemize}[leftmargin=0.5cm, nosep]
        \item Développement d'une application inspirée de Waze, spécifiquement adaptée aux missions de sécurité publique (React Native, Flask).
    \end{itemize}
    \vspace{0.2cm}
    \item \textbf{Étude Spatio-Statistique (Île-de-France)} :
    \begin{itemize}[leftmargin=0.5cm, nosep]
        \item Modélisation de l'impact des modifications du réseau de transport sur le report modal domicile-travail via une approche panel spatio-temporelle et simulation de scénarios prospectifs (Scikit-Learn, Seaborn, SciPy, Proj, GeoPandas).
    \end{itemize}
    \vspace{0.2cm}
    \item \textbf{Plateforme SaaS KlassIvoire} :
    \begin{itemize}[leftmargin=0.5cm, nosep]
        \item Contribution au développement pour le suivi scolaire numérique, avec intégration d’outils d'IA pour l'analyse des performances.
        \item Technologies : Django, SQL, PostgreSQL, Docker, Cloudflare R2, Redis, API Gemini.
    \end{itemize}
    \vspace{0.2cm}
    \item \textbf{Plateforme d'analyse d'établissements scolaires} :
    \begin{itemize}[leftmargin=0.5cm, nosep]
        \item Conception d’une plateforme intelligente de suivi géographique permettant l’exploitation de données spatiales pour la prise de décision.
        \item Technologies : Django, PostgreSQL/PostGIS, Leaflet.js.
    \end{itemize}
\end{itemize}


% ================= VIE ASSOCIATIVE =================
\section*{Vie Associative}

\cventry{Professeur particulier}{Île-de-France}{Cours à domicile -- Mathématiques et Physique}{Depuis Oct. 2023}
\noindent Pédagogie et méthodologie, écoute attentive des besoins des élèves.

\cventry{VERTIGEO -- Junior-Entreprise de l'ENSG Géomatique}{Champs-sur-Marne}{Secrétaire Général}{Fév. 2024 -- Mars 2025}
\noindent Comptes rendus, archivage, gestion des relations avec les partenaires administratifs, organisation de réunions.

\cventry{Épicerie solidaire campésienne}{Champs-sur-Marne}{Bénévole}{Fév. 2024}
\noindent Collecte de produits alimentaires et d'hygiène dans un magasin Carrefour. Promotion de l'échange et de l'ouverture d'esprit.

% ================= LANGUES ET INTÉRÊTS =================
\section*{Langues et Centres d'intérêt}
\noindent \textbf{Langues :} Français (Langue maternelle), Anglais (TOEIC - B2) \\[0.1cm]
\textbf{Centres d'intérêt :} Intelligence Artificielle, Veille technologique, Football, Gymnastique acrobatique, Jeux de société

\end{document}